\begin{comment}
    Introducción
        XAI y las distintas métricas que existen
            Feature attribution, counterfactual, rule based, sufficient reason, etc
        Algoritmos de feature attribution
        Código fuente 
\end{comment}

Las capacidades de los distintos modelos de inteligencia artificial (IA) han ido creciendo exponencialmente en los últimos años. Tareas que creíamos imposibles de realizar a través de estas técnicas, como la resolución de problemas matemáticos complejos, hoy en día pueden ser resueltas por los LLMs \cite{frontierMath}. No sólo aumentó la complejidad de las tareas que pueden resolver, sino también la complejidad de sus arquitecturas y de como están compuestos. Los modelos de hace unos años tenían 1e8 parámetros, mientras que hoy en día ya llegan a 1e12. \cite{EpochLargeScaleModels2024}. El problema que esto genera es que cada vez es más difícil entender como funcionan y como obtienen las respuestas que nos dan. Esto es algo sumamente importante, ya que si uno está utilizando estas técnicas para diagnósticos médicos u otras situaciones igual de delicadas, es fundamental  entender como los modelos alcanzan sus conclusiones . Ahí es donde entra en juego el área de XAI, Explainable Artificial Intelligence.

En esta tesis nos centramos en \textit{SHAP} \cite{shapOriginalPaper}, el cual es un método que, en principio, provee explicaciones agnósticas al modelo, con un alcance local y que es del tipo de feature attribution. Esto significa que para una predicción individual que realiza el modelo, \textit{SHAP} va a devolver un valor asociado a cada feature de la instancia. Decimos que es agnóstico, puesto que su valor depende meramente del ouput del modelo, por lo que no esta acoplado a su representación interna. Aún así, \textit{SHAP} actualmente es un framework que engloba varios tipos de métodos, ya que hay distintas aproximaciones según cuál sea tu tipo de modelo. Por ejemplo, existen métodos como DeepExplainer, TreeExplainer y LinearExplainer dentro del  \textcolor{blue}{\href{https://shap.readthedocs.io/en/latest/index.html}{framework}}  SHAP, los cuales se utilizan para modelos específicos, pero también hay otros métodos como KernelExplainer, que son agnósticos al modelo.

%\santi{Pero no es este el motivo por el cual decimos que es agnóstico, no? Decimos que es agnóstico porque su valor depende únicamente del mappeo de entradas a salidas. O sea, es una propiedad del mappeo subyacente al modelo, ignorando la arquitectura específica. Entiendo que la idea original es esa, y estas versiones que se nombras después son aproximaciones imperfectas (pq muchas veces nisiquiera coinciden con los shapley o SHAP) para modelos particulares.}

Uno de los problemas principales de este framework es lo costoso que es calcularlo, y en particular, hay estudios que analizan esta dificultad desde el punto de vista de la complejidad computacional. Por ejemplo, se demostró que calcular SHAP para un clasificador trivial era NP-completo, o intratable, para datos con distribuciones más o igual de complejas que una Naive Bayes \cite{arenas2021tractability}.
A raíz de esto, nos resultó interesante analizar si existía alguna variación de SHAP, la cual pueda ser calculada en tiempo polinomial para una distribución de red bayesiana. Además, actualmente hay trabajos similares a este, que buscan analizar la complejidad de SHAP según la variante, el modelo y la distribución \cite{marzouk2025computationaltractabilitymanyshapley} %\santi{Hace poco salió este paper, habría que agregarlo en algún lado -> On the Computational Tractability of the (Many) Shapley Values} 

SHAP está basado en los Shapley values \cite{shapley1953value}, un valor de teoría de juegos.  Una variación de los mismos son los assymetric shapley values (ASV), los cuales pueden introducir el factor de causalidad al cálculo de los mismos. De esta manera, se define ASV \cite{frye2019asymmetric}, un framework similar a SHAP, el cual además tiene en cuenta el grafo causal de los datos. El objetivo de esta tesis es lograr calcular ASV en tiempo polinomial de forma exacta y aproximada para datos con distribuciones de redes bayesianas.

\subsection{Código fuente}

Todos los algoritmos presentados en este trabajo fueron implementados. El código fuente puede ser encontrado online en el siguiente repositorio: 

\begin{table}[H]
\centering
\begin{tabular}{ll}
\toprule
\textbf{Repository} & \textbf{Repository} \\
\midrule
Source code & \url{https://github.com/EchuCompa/pasantia-BICC} \\
\end{tabular}
\end{table}